% !TEX root = ../../../main.tex
% 所属：第四部分：Fourier 分析 / 第4章：$\R$上的Fourier变换
% 节：4.2 Schwartz 空间
% 正文按原文件顺序保存；由 ../chapter4.tex 调用。
\section{Schwartz 空间}

\begin{definition}
\(\R\) 上的 \textbf{Schwartz 空间}由满足下列条件的函数 \(f\) 组成：\(f\) 无穷可导，并且对于任意非负整数 \(k,l\)，都有
\[
\sup_{x\in\R}\abs{x}^{k}\abs{f^{(l)}(x)}<\infty.
\]
这意味着 \(f\) 及其各阶导数都在无穷远处速降。记这一空间为
\[
\mathcal S=\mathcal S(\R).
\]
在通常的函数加法与数乘下，\(\mathcal S(\R)\) 是 \(\C\) 上的向量空间。
\end{definition}

由定义可知，若 \(f\in\mathcal S(\R)\)，则
\[
f'(x)=\frac{\dd f}{\dd x}(x)\in\mathcal S(\R),
\qquad
xf(x)\in\mathcal S(\R).
\]
更一般地，Schwartz 空间关于求导以及与多项式相乘都是封闭的。

\begin{definition}
若 \(f\in\mathcal S(\R)\)，则其 Fourier 变换定义为
\[
\widehat f(\xi)
=\int_{-\infty}^{\infty}f(x)e^{-2\pi i x\xi}\dd x.
\]
以下用记号
\[
f(x)\longmapsto\widehat f(\xi)
\]
表示 \(\widehat f\) 是 \(f\) 的 Fourier 变换。
\end{definition}

\begin{proposition}\label{prop:fourier-basic-properties-schwartz}
若 \(f\in\mathcal S(\R)\)，则 Fourier 变换满足下列性质：
\begin{enumerate}[label=(\roman*)]
\item 对任意 \(h\in\R\)，
\[
f(x+h)\longmapsto e^{2\pi i h\xi}\widehat f(\xi).
\]

\item 对任意 \(h\in\R\)，
\[
f(x)e^{-2\pi i xh}\longmapsto\widehat f(\xi+h).
\]

\item 若 \(\delta>0\)，则
\[
f(\delta x)\longmapsto\delta^{-1}\widehat f(\delta^{-1}\xi).
\]

\item
\[
f'(x)\longmapsto2\pi i\xi\widehat f(\xi).
\]

\item
\[
-2\pi ixf(x)\longmapsto\frac{\dd}{\dd\xi}\widehat f(\xi).
\]
\end{enumerate}
\end{proposition}

\begin{proof}
性质 (i)--(iii) 分别由变量替换、定义以及伸缩性直接得到。

对 (iv)，在 \([-N,N]\) 上分部积分可得
\[
\int_{-N}^{N}f'(x)e^{-2\pi i x\xi}\dd x
=\left[f(x)e^{-2\pi i x\xi}\right]_{-N}^{N}
+2\pi i\xi\int_{-N}^{N}f(x)e^{-2\pi i x\xi}\dd x.
\]
由于 \(f\) 速降，令 \(N\to\infty\) 即得结论。

最后，由
\[
\frac{\widehat f(\xi+h)-\widehat f(\xi)}{h}
=\int_{-\infty}^{\infty}
 f(x)e^{-2\pi i x\xi}
 \frac{e^{-2\pi i xh}-1}{h}\dd x
\]
以及
\[
\abs{\frac{e^{-2\pi i xh}-1}{h}}\leq2\pi\abs{x},
\]
而 \(xf(x)\) 可积，故由控制收敛定理可将极限移入积分号，得到
\[
\frac{\dd}{\dd\xi}\widehat f(\xi)
=\int_{-\infty}^{\infty}(-2\pi ix)f(x)e^{-2\pi i x\xi}\dd x.
\]
这就证明了 (v)。
\end{proof}

\begin{theorem}\label{thm:fourier-preserves-schwartz}
若 \(f\in\mathcal S(\R)\)，则
\[
\widehat f\in\mathcal S(\R).
\]
\end{theorem}

\begin{proof}
对任意非负整数 \(k,l\)，由命题 \ref{prop:fourier-basic-properties-schwartz} 可知
\[
\xi^k\left(\frac{\dd}{\dd\xi}\right)^l\widehat f(\xi)
\]
是函数
\[
\frac{1}{(2\pi i)^k}
\left(\frac{\dd}{\dd x}\right)^k
\left[(-2\pi ix)^l f(x)\right]
\]
的 Fourier 变换。括号中的函数属于 \(\mathcal S(\R)\)，故其 Fourier 变换有界。因此
\[
\sup_{\xi\in\R}
\abs{\xi}^{k}
\abs{\widehat f^{(l)}(\xi)}<\infty,
\]
从而 \(\widehat f\in\mathcal S(\R)\)。
\end{proof}

下一节将通过仔细研究 Gaussian 函数 \(e^{-ax^2}\) 来证明 Fourier 变换的反演公式，因此我们提前给出他的一些性质。
容易看出，当 \(a>0\) 时，函数 \(e^{-ax^2}\) 属于 \(\mathcal S(\R)\)。

\begin{proposition}\label{prop:gaussian-normalization}
取 \(a=\pi\)，则有标准化公式
\[
\int_{-\infty}^{\infty}e^{-\pi x^2}\dd x=1.
\]
\end{proposition}

\begin{proof}
记
\[
I=\int_{-\infty}^{\infty}e^{-\pi x^2}\dd x>0.
\]
由极坐标变换，
\[
\begin{aligned}
I^2
&=\int_{-\infty}^{\infty}\int_{-\infty}^{\infty}
  e^{-\pi(x^2+y^2)}\dd x\dd y\\
&=\int_0^{2\pi}\int_0^{\infty}e^{-\pi r^2}r\dd r\dd\theta
 =\int_0^{\infty}2\pi r e^{-\pi r^2}\dd r\\
&=\left[-e^{-\pi r^2}\right]_{0}^{\infty}=1.
\end{aligned}
\]
因此 \(I=1\)。
\end{proof}

\begin{proposition}\label{thm:gaussian-fourier-transform}
若$f(x)=e^{-\pi x^2},$则
\[
\widehat f(\xi)=f(\xi)=e^{-\pi\xi^2}.
\]
\end{proposition}

\begin{proof}
令
\[
F(\xi)=\widehat f(\xi)
=\int_{-\infty}^{\infty}e^{-\pi x^2}e^{-2\pi i x\xi}\dd x.
\]
由命题 \ref{prop:gaussian-normalization}，\(F(0)=1\)。又因
\[
f'(x)=-2\pi xf(x),
\]
结合命题 \ref{prop:fourier-basic-properties-schwartz} 的 (iv)、(v)，有
\[
F'(\xi)
=i\int_{-\infty}^{\infty}f'(x)e^{-2\pi i x\xi}\dd x
=i(2\pi i\xi)F(\xi)
=-2\pi\xi F(\xi).
\]
于是
\[
\frac{\dd}{\dd\xi}\left(F(\xi)e^{\pi\xi^2}\right)=0.
\]
由 \(F(0)=1\) 可知
\[
F(\xi)e^{\pi\xi^2}=1,
\]
故 \(F(\xi)=e^{-\pi\xi^2}=\widehat f(\xi)\)。
\end{proof}

\begin{corollary}\label{cor:scaled-gaussian-fourier-transform}
若 \(\delta>0\)，且
\[
K_\delta(x)=\delta^{-1/2}e^{-\pi x^2/\delta},
\]
则
\[
\widehat{K_\delta}(\xi)=e^{-\pi\delta\xi^2}.
\]
\end{corollary}

\begin{proof}
将 \(K_\delta\) 写成
\[
K_\delta(x)=\delta^{-1/2}f(\delta^{-1/2}x),
\qquad f(x)=e^{-\pi x^2},
\]
再应用 \ref{prop:fourier-basic-properties-schwartz} 的 (iii) 与 \ref{thm:gaussian-fourier-transform} 即得。
\end{proof}

\begin{theorem}\label{thm:gaussian-good-kernel}
当 \(\delta\to0^+\) 时，函数族 \(\{K_\delta\}_{\delta>0}\) 是一族好核。也就是说，对于
\[
K_\delta(x)=\delta^{-1/2}e^{-\pi x^2/\delta},
\]
有：
\begin{enumerate}[label=(\roman*)]
\item
\[
\int_{-\infty}^{\infty}K_\delta(x)\dd x=1;
\]

\item 存在与 \(\delta\) 无关的常数 \(M\)，使得
\[
\int_{-\infty}^{\infty}\abs{K_\delta(x)}\dd x\leq M;
\]

\item 对任意 \(\eta>0\)，当 \(\delta\to0^+\) 时，
\[
\int_{\abs{x}>\eta}\abs{K_\delta(x)}\dd x\longrightarrow0.
\]
\end{enumerate}
\end{theorem}

\begin{proof}
在 (i) 中令 \(x=\delta^{1/2}y\)，由命题 \ref{prop:gaussian-normalization} 得
\[
\int_{-\infty}^{\infty}K_\delta(x)\dd x
=\int_{-\infty}^{\infty}e^{-\pi y^2}\dd y=1.
\]
由于 \(K_\delta\geq0\)，故 (ii) 可取 \(M=1\)。最后，同样作变量替换可得
\[
\int_{\abs{x}>\eta}\abs{K_\delta(x)}\dd x
=\int_{\abs{y}>\eta/\delta^{1/2}}e^{-\pi y^2}\dd y
\longrightarrow0,
\]
这就证明了 (iii)。
\end{proof}

\begin{definition}
对 \(f,g\in\mathcal S(\R)\)，定义它们的\textbf{卷积}为
\[
(f*g)(x)
=\int_{-\infty}^{\infty}f(x-t)g(t)\dd t.
\]
由于 \(f,g\) 都是速降函数，上述积分绝对收敛，因而定义良好。
\end{definition}

\begin{proposition}\label{prop:gaussian-approximate-identity}
若 \(f\in\mathcal S(\R)\)，则当 \(\delta\to0^+\) 时，关于 \(x\in\R\) 一致地有
\[
(f*K_\delta)(x)\longrightarrow f(x).
\]
\end{proposition}

\begin{proof}
首先，断言\(f\) 在 \(\R\) 上一致连续。事实上，给定 \(\eps>0\)，可取 \(R>0\)，使得 \(\abs{x}>R\) 时 \(\abs{f(x)}<\eps/4\)；而 \(f\) 在紧区间 \([-R-1,R+1]\) 上一致连续。两者结合可知，存在 \(\eta>0\)，使得当 \(\abs{x-y}<\eta\) 时，
\[
\abs{f(x)-f(y)}<\eps.
\]

由定理 \ref{thm:gaussian-good-kernel} 的 (i)，
\[
(f*K_\delta)(x)-f(x)
=\int_{-\infty}^{\infty}K_\delta(t)
\bigl(f(x-t)-f(x)\bigr)\dd t.
\]
于是
\[
\begin{aligned}
\abs{(f*K_\delta)(x)-f(x)}
\leq{}&\int_{\abs{t}>\eta}K_\delta(t)
\abs{f(x-t)-f(x)}\dd t\\
&+\int_{\abs{t}\leq\eta}K_\delta(t)
\abs{f(x-t)-f(x)}\dd t.
\end{aligned}
\]
记 \(B=\sup_{x\in\R}\abs{f(x)}\)。第一项不超过
\[
2B\int_{\abs{t}>\eta}K_\delta(t)\dd t,
\]
由好核的性质 (iii)，当 \(\delta\to0^+\) 时它一致趋于 \(0\)；第二项不超过
\[
\eps\int_{-\infty}^{\infty}K_\delta(t)\dd t=\eps.
\]
因此 \((f*K_\delta)(x)\to f(x)\) 关于 \(x\) 一致成立。
\end{proof}
